\textbf{Phần 1. Trắc nghiệm lựa chọn}
\setcounterpageref{ex}{0}

\begin{ex} %[2D2N09]
Trên cành cây có 8 con chim, có thêm 4 con chim bay đến. Hỏi trên cành cây có tất cả bao nhiêu con chim?
\choice
{$11$ con}
{$13$ con}
{$14$ con}
{\True $12$ con}
\loigiai{
$8 + 4 = 12$ con chim.\\
\textbf{Đáp án của câu hỏi là D}
}
\end{ex}

\begin{ex} %[2D2N09]
Trong bể có 15 con cá, vớt ra 6 con cá. Trong bể còn lại số con cá là
\choice
{\True $9$ con}
{$8$ con}
{$10$ con}
{$11$ con}
\loigiai{
$15 - 6 = 9$ con cá.\\
\textbf{Đáp án của câu hỏi là A}
}
\end{ex}

\begin{ex} %[2D2H09]
Mai có 9 cái kẹo, mẹ cho Mai thêm 5 cái kẹo nữa. Mai có tất cả số cái kẹo là
\choice
{$13$ cái}
{$15$ cái}
{\True $14$ cái}
{$16$ cái}
\loigiai{
$9 + 5 = 14$ cái kẹo.\\
\textbf{Đáp án của câu hỏi là C}
}
\end{ex}

\begin{ex} %[2D2H09]
Lan có 14 quả bóng bay, Lan làm vỡ 5 quả. Lan còn lại số quả bóng bay là
\choice
{$8$ quả}
{\True $9$ quả}
{$10$ quả}
{$11$ quả}
\loigiai{
$14 - 5 = 9$ quả.\\
\textbf{Đáp án của câu hỏi là B}
}
\end{ex}

\begin{ex} %[2D2H09]
Một xe buýt có 12 hành khách, đến bến đỗ có thêm 6 hành khách lên xe. Trên xe có tất cả số hành khách là
\choice
{$17$ người}
{$19$ người}
{\True $18$ người}
{$20$ người}
\loigiai{
$12 + 6 = 18$ người.\\
\textbf{Đáp án của câu hỏi là C}
}
\end{ex}

\begin{ex} %[2D2H09]
Một bãi đỗ xe có 16 chiếc ô tô, sau đó có 7 chiếc ô tô rời bãi. Trong bãi còn lại số ô tô là
\choice
{\True $9$ chiếc}
{$8$ chiếc}
{$10$ chiếc}
{$11$ chiếc}
\loigiai{
$16 - 7 = 9$ chiếc.\\
\textbf{Đáp án của câu hỏi là A}
}
\end{ex}

\begin{ex} %[2D2V09]
Hà có 8 quyển vở, Hà được thưởng thêm 7 quyển vở nữa. Hỏi Hà có tất cả bao nhiêu quyển vở?
\choice
{$14$ quyển}
{\True $15$ quyển}
{$16$ quyển}
{$17$ quyển}
\loigiai{
$8 + 7 = 15$ quyển vở.\\
\textbf{Đáp án của câu hỏi là B}
}
\end{ex}

\begin{ex} %[2D2V09]
Đàn gà có 17 con, đã bán đi 8 con. Hỏi đàn gà còn lại bao nhiêu con?
\choice
{$8$ con}
{$10$ con}
{$11$ con}
{\True $9$ con}
\loigiai{
$17 - 8 = 9$ con.\\
\textbf{Đáp án của câu hỏi là D}
}
\end{ex}

\begin{ex} %[2D2V09]
Hộp thứ nhất có 9 cái bút, hộp thứ hai có nhiều hơn hộp thứ nhất 4 cái bút. Hỏi hộp thứ hai có bao nhiêu cái bút?
\choice
{\True $13$ cái}
{$12$ cái}
{$14$ cái}
{$15$ cái}
\loigiai{
$9 + 4 = 13$ cái bút.\\
\textbf{Đáp án của câu hỏi là A}
}
\end{ex}

\begin{ex} %[2D2V09]
Cửa hàng có 13 thùng sữa, đã bán đi 6 thùng. Cửa hàng còn lại số thùng sữa là
\choice
{$6$ thùng}
{$8$ thùng}
{$9$ thùng}
{\True $7$ thùng}
\loigiai{
$13 - 6 = 7$ thùng.\\
\textbf{Đáp án của câu hỏi là D}
}
\end{ex}

\begin{ex} %[2D2V09]
Trên cành có 9 quả táo chín và 8 quả táo xanh. Trên cành có tất cả số quả táo là
\choice
{$16$ quả}
{$18$ quả}
{\True $17$ quả}
{$19$ quả}
\loigiai{
$9 + 8 = 17$ quả.\\
\textbf{Đáp án của câu hỏi là C}
}
\end{ex}

\begin{ex} %[2D2C09]
Một xe khách có 15 hành khách. Đến một bến đỗ, có 4 người xuống xe và 6 người lên xe. Hỏi trên xe lúc này có bao nhiêu hành khách?
\choice
{$16$ người}
{\True $17$ người}
{$18$ người}
{$19$ người}
\loigiai{
$15 - 4 + 6 = 17$ người.\\
\textbf{Đáp án của câu hỏi là B}
}
\end{ex}

\textbf{Phần 2. Trắc nghiệm đúng sai}
\setcounterpageref{ex}{0}

\begin{ex} %[2D2H09]
Bác An có 9 con vịt, bác mua thêm 6 con vịt nữa. Xét tính đúng sai:
\choiceTF
{\True Bác An mua thêm nghĩa là ta làm phép tính cộng.}
{\True Bác An có tất cả 15 con vịt ($9 + 6 = 15$).}
{Nếu bác An bán bớt 5 con vịt thì còn lại 11 con.}
{\True Khi đó bác An còn lại $15 - 5 = 10$ con vịt.}
\loigiai{
\begin{itemize}
    \item Mua thêm là phép cộng: Đúng.
    \item $9 + 6 = 15$: Đúng.
    \item $15 - 5 = 10 \neq 11$. Ý 3 Sai, Ý 4 Đúng.
\end{itemize}
\begin{itemize}
    \item a) \textbf{Đúng}.
    \item b) \textbf{Đúng}.
    \item c) \textbf{Sai}.
    \item d) \textbf{Đúng}.
\end{itemize}
}
\end{ex}

\begin{ex} %[2D2V09]
Một đĩa có 16 quả dâu tây, bé Nam ăn mất 7 quả. Xét tính đúng sai:
\choiceTF
{\True Số dâu tây còn lại là 9 quả ($16 - 7 = 9$).}
{\True Bé Nam ăn mất là bài toán về bớt một số đơn vị.}
{Nếu mẹ cho thêm vào đĩa 4 quả thì trong đĩa có 12 quả.}
{\True Khi đó trong đĩa có $9 + 4 = 13$ quả dâu tây.}
\loigiai{
\begin{itemize}
    \item Còn lại: $16 - 7 = 9$: Đúng.
    \item Bài toán bớt đi: Đúng.
    \item $9 + 4 = 13 \neq 12$. Ý 3 Sai, Ý 4 Đúng.
\end{itemize}
\begin{itemize}
    \item a) \textbf{Đúng}.
    \item b) \textbf{Đúng}.
    \item c) \textbf{Sai}.
    \item d) \textbf{Đúng}.
\end{itemize}
}
\end{ex}

\textbf{Phần 3. Trả lời ngắn}
\setcounterpageref{ex}{0}

\begin{ex} %[2D2H09]
Có 8 con bướm, thêm 5 con bướm nữa bay tới. Hỏi có tất cả bao nhiêu con bướm?
\shortans{13}
\loigiai{
\textbf{Đáp số:} $13$\\
$8 + 5 = 13$ con.
}
\end{ex}

\begin{ex} %[2D2H09]
Có 15 bông hoa, cắm vào lọ 7 bông hoa. Còn lại bao nhiêu bông hoa chưa cắm?
\shortans{8}
\loigiai{
\textbf{Đáp số:} $8$\\
$15 - 7 = 8$ bông.
}
\end{ex}

\begin{ex} %[2D2V09]
Lan có 9 cái kẹo, mẹ cho thêm 7 cái kẹo. Hỏi Lan có tất cả bao nhiêu cái kẹo?
\shortans{16}
\loigiai{
\textbf{Đáp số:} $16$\\
$9 + 7 = 16$ cái.
}
\end{ex}

\begin{ex} %[2D2V09]
Một tổ có 14 bạn, có 6 bạn đi làm nhiệm vụ khác. Hỏi tổ đó còn lại bao nhiêu bạn?
\shortans{8}
\loigiai{
\textbf{Đáp số:} $8$\\
$14 - 6 = 8$ bạn.
}
\end{ex}

\begin{ex} %[2D2V09]
Trong chuồng có 8 con gà trống và 8 con gà mái. Hỏi trong chuồng có tất cả bao nhiêu con gà?
\shortans{16}
\loigiai{
\textbf{Đáp số:} $16$\\
$8 + 8 = 16$ con.
}
\end{ex}

\begin{ex} %[2D2C09]
Nam có 12 viên bi. Nam cho Bình 3 viên bi và được Dũng cho thêm 5 viên bi. Hỏi Nam có bao nhiêu viên bi?
\shortans{14}
\loigiai{
\textbf{Đáp số:} $14$\\
$12 - 3 + 5 = 14$ viên bi.
}
\end{ex}
